\documentclass[10pt,a4paper]{article}
\usepackage[margin=20mm]{geometry}
\usepackage{amsmath,amssymb,booktabs,array,hyperref}
\hypersetup{colorlinks=true,urlcolor=blue}
\setlength{\parindent}{0pt}
\setlength{\parskip}{5pt}
\newcommand{\R}{\mathbb R}
\newcommand{\C}{\mathbb C}
\title{A unit Lipschitz field whose full Hilbert transform\newline has no square-integrable principal-value representative}
\author{lizaixi01 \qquad TLMC 00000001911}
\date{24 October 2026}
\begin{document}
\maketitle
\vspace{-6mm}

\textbf{Scope.} The printed first conjunct says that
$H_Vf(x)=\operatorname{p.v.}\int f(x-tV(x))\,dt/t$
is $L^2$ bounded for Lipschitz vector fields. No outer cutoff is printed.
We disprove this \emph{full, untruncated strong-$L^2$ assertion} on the usual
Euclidean plane, even for an everywhere unit, globally Lipschitz field and
a genuine Schwartz input. This is not a counterexample to the conventional
short-scale Stein question, its weak-type formulation, or the source's
undefined resonance clause. Lacey--Li explicitly use a finite outer scale
in their initial definition and in the Stein conjecture (printed pp.~1--2);
their printed p.~27 separately discusses the full transform [1].

\textbf{1. Actual field and test function.}
Identify the Euclidean plane with $\C$, with its native norm and Lebesgue measure.
For $z=x+iy$, put
\[
 W(z)=\max(1,x)+iy,\qquad V(z)=\|W(z)\|^{-1}W(z).
\]
The real clamp is $1$-Lipschitz. Squaring the Euclidean norms shows that
$W$ is $1$-Lipschitz, and $\|W(z)\|\ge1$. For $\|u\|,\|v\|\ge1$,
\[
 \left\|\frac{u}{\|u\|}-\frac{v}{\|v\|}\right\|
 \le \frac{\|u-v\|+|\|v\|-\|u\||}{\|u\|}
 \le2\|u-v\|.
\]
Consequently $V$ is globally $2$-Lipschitz, has norm one everywhere,
and equals $z/\|z\|$ when $x\ge1$.
Choose the native smooth bump $f$ centered at zero with inner radius $1$
and outer radius $2$. It satisfies $0\le f\le1$, $f=1$ on the closed unit ball,
and $f=0$ for $\|z\|\ge2$. Its actual compact support and all-order smoothness
construct a native Schwartz map, hence $f\in L^2(\C)$.

\textbf{2. Genuine principal values on a cone.}
For real $\varepsilon,R$, define the actual cutoff integral
\[
 H_{\varepsilon,R}f(z)=
 \int_{\{t:\,\varepsilon\le |t|\le R\}}
       \frac{f(z-tV(z))}{t}\,dt.
\]
Let $E=\{x+iy:x\ge4,\ 0\le y\le x\}$ and $r=\|z\|$.
For $z\in E$, one has $4\le r\le2x$ and
$\|z-tV(z)\|=|r-t|$. Thus the actual kernel vanishes outside
$[r-2,r+2]\subset(0,\infty)$, is nonnegative everywhere, and is continuous
on that interval. Its global integral is genuinely integrable, and every
cutoff restriction is integrable. Whenever $\varepsilon\le r-2$ and $R\ge r+2$,
\[
 H_{\varepsilon,R}f(z)=P(z):=
 \int_{r-2}^{r+2}\frac{f(z-tV(z))}{t}\,dt.
\]
This eventual equality proves the real two-parameter limit
$H_{\varepsilon,R}f(z)\to P(z)$ as
$\varepsilon\to0^+$ and $R\to+\infty$ in the product filter.
No value is assigned by a totalized limit outside this cone.

\textbf{3. A lower bound and the $L^2$ obstruction.}
On $[r-1,r+1]$, the bump equals one. Positivity and interval monotonicity give
\[
 P(z)\ge\int_{r-1}^{r+1}\frac{dt}{t}
 \ge\frac{2}{r+1}\ge\frac1r\ge\frac1{2x}\quad(z\in E).
\]
Suppose a native $L^2$ function $g$ agrees almost everywhere on $E$ with
these actual principal-value limits. Uniqueness of limits implies
$g(z)=P(z)$ there. Native $L^2$ membership makes $g^2$ integrable.
Transport through the actual volume-preserving real-coordinate measurable
equivalence, followed by Fubini, gives an integrable row integral
$G(x)=\int_\R g(x+iy)^2\,dy$. For almost every $x\ge4$,
\[
 G(x)\ge\int_0^x\frac{dy}{4x^2}=\frac1{4x}.
\]
Integrable domination would make $x\mapsto1/x$ integrable on $[4,\infty)$,
contrary to native nonintegrability of that function. Hence no such $g$ exists.

\newpage
\textbf{4. Exact operator consequence.}
The formal bounded-operator claim uses the native Banach space
$L^2(\C;\R)$ and an actual continuous real-linear map
$T:L^2(\C;\R)\to L^2(\C;\R)$. Its defining compatibility is: for every
native real-valued Schwartz map $h$, for almost every $z$,
\[
 H_{\varepsilon,R}h(z)\longrightarrow (T h)(z)
 \quad\text{as }(\varepsilon,R)\to(0^+,+\infty).
\]
This is a necessary compatibility for an $L^2$-bounded linear extension
of the full principal-value transform. Applying it to the constructed
$f$ would provide the impossible $L^2$ representative above. The result
therefore excludes every such bounded operator. The stronger intermediate
result excludes even an arbitrary $L^2$ representative on $E$; linearity
or an operator-norm estimate is not needed for that contradiction.
The proof does not posit the full transform's values elsewhere.

\textbf{5. Complete formal coverage and interpretation.}
The pinned Lean project proves the construction directly. It does not add
an axiom asserting a literature theorem, a desired transform value, or a
divergent norm. The geometric and measure identifications are native:
$\|z\|=\sqrt{x^2+y^2}$ and the real-coordinate map preserves native Lebesgue
volume. In particular, the product measure used for Fubini does not replace
the Euclidean metric by the product sup metric. The restriction to real
Schwartz tests already rules out the printed real strong-$L^2$ transform.

\small
\begin{tabular}{@{}p{.28\linewidth}p{.67\linewidth}@{}}
\toprule
Part & Principal declarations in namespace \texttt{TLMC1911} \\
\midrule
Euclidean field & \texttt{euclidean\_norm}, \texttt{euclidean\_volume},
\texttt{clipped\_lipschitz}, \texttt{normalization\_bound},
\texttt{field\_unit}, \texttt{field\_lipschitz}, \texttt{field\_radial} \\
Smooth input & \texttt{bump}, \texttt{test}, \texttt{test\_one},
\texttt{test\_zero}, \texttt{test\_memLp} \\
Actual cutoff and PV & \texttt{radial\_line}, \texttt{kernel\_integrable},
\texttt{cutoff\_integrable}, \texttt{kernel\_plateau},
\texttt{cutoff\_eq\_conePV}, \texttt{conePV\_tendsto}, \texttt{conePV\_lower} \\
Native $L^2$ contradiction & \texttt{no\_memLp\_cone\_lower},
\texttt{no\_L2\_PV\_representative}, \texttt{RepresentsFullPV},
\texttt{no\_bounded\_fullPV\_operator},
\texttt{printed\_untruncated\_conjunct\_false} \\
\bottomrule
\end{tabular}
\normalsize

The final theorem exhibits a genuine $2$-Lipschitz unit field for which
there is no native continuous-linear $L^2$ operator satisfying this full-PV
compatibility. All variables in the cutoff limits, all almost-everywhere
statements, and the square-integrability assertions retain their actual
native meanings. The report uses the conventional planar interpretation,
which the printed text does not further specify. It only refutes the first
untruncated strong-$L^2$ conjunct. The conventional outer-truncated,
short-scale Stein problem and the resonance-set statement remain outside
the claim. A complete copy of both original language versions is in
\texttt{SOURCE.md}.

\textbf{6. Reproduction and evidence.}
Lean is pinned to 4.33.0, with Mathlib revision
\begin{quote}\ttfamily
db584cd6d46c92f209a44c0f1c829460d327499d
\end{quote}
From \texttt{lean/}, run:
\begin{quote}\ttfamily
lean --version\\
lake build +Main\\
lake env lean -DwarningAsError=true Check.lean\\
lake env leanchecker --fresh Main
\end{quote}
Alternatively run \texttt{lean --version}, then \texttt{bash verify.sh} for the
remaining three commands. The audit invokes ordinary \texttt{\#print axioms}
for every actual owned declaration, including generated/private names,
and checks exact namespace coverage. Only \texttt{propext},
\texttt{Classical.choice}, and \texttt{Quot.sound} are permitted.
The package retains the full source-bound command logs, execution and
CPU-affinity samples, qualification and PDF/semantic reviews.
No production timing is offered as causal research-efficiency evidence.

\textbf{Reference.}
[1] M. Lacey and X. Li, \emph{On a Conjecture of E. M. Stein on the Hilbert
Transform on Vector Fields}, arXiv:0704.0808v3 (11 January 2008).
\url{https://arxiv.org/abs/0704.0808}.
The complete selected printed pp.~1--2 and 27, including their page images,
were read; no claim is made to have read the whole 94-page work.
They distinguish the usual finite outer scale from the full transform;
the counterexample and every formal inference above are proved directly.
\end{document}
